\chapter{Proofs: Sufficiency Relative to a Closed Task and System}
\label{ch:proofs}

\begin{quotation}
\noindent\textit{Synopsis.} This chapter addresses the strongest apparent
counterexample to the monograph's thesis: a formal proof, once written,
appears complete. The chapter shows this is correct precisely for
verification relative to a fixed axiom system, and precisely why it stops
being correct the moment explanation, intuition, generalization, pedagogy,
or future theoretical change become the task instead. Proofs therefore
demonstrate that task-relative sufficiency is genuinely achievable in the
closed case, while simultaneously reinforcing why task boundaries matter as
much as they do throughout the rest of the book.
\end{quotation}

\section{The Strongest Apparent Counterexample}

Any framework claiming that no finite substrate is generally sufficient
invites the objection: but a machine-checkable proof, once written, is
complete for verification. This chapter states precisely why that
objection is correct as far as it goes, and precisely where it stops
applying -- a role this chapter is written to make explicit, since the
objection would otherwise hang unanswered over the rest of the
monograph.

\begin{proposition}[Verification Sufficiency]
\label{prop:verification-sufficiency}
For a fixed axiom system $\mathcal{S}$, a machine-checkable proof $p$ is
sufficient (\cref{def:sufficiency}) for the task of verification relative
to $\mathcal{S}$.
\end{proposition}

\begin{proofsketch}
The verification algorithm for $\mathcal{S}$ is finite and decidable by
construction; every distinction it requires to certify $p$ appears
explicitly in $p$ together with $\mathcal{S}$'s axioms. No distinction
required by the verification task is discarded, so
\cref{def:sufficiency} is satisfied.
\end{proofsketch}

\begin{corollary}[Explanation Is Not Guaranteed]
\label{cor:explanation-not-guaranteed}
\cref{prop:verification-sufficiency} does not extend to the tasks of
explanation, generalization, or adaptation to a changed axiom system.
\end{corollary}

\begin{proofsketch}
These tasks lie outside the closed task space $\{\tau_{\text{verify}}\}$
for which $\mathcal{S}$ and $p$ were held fixed in
\cref{prop:verification-sufficiency}'s hypothesis; by
\cref{prop:task-unavailability}, sufficiency for tasks outside a closed,
anticipated set is not guaranteed by a finite substrate.
\end{proofsketch}

\section{Why This Resolves Rather Than Merely Sidesteps the Objection}

\cref{prop:verification-sufficiency} shows the impossibility claim of
\cref{principle:global-insufficiency} was never intended to apply where
both the task and the formal system are closed and fixed in advance;
\cref{cor:explanation-not-guaranteed} shows that the moment either is
allowed to vary, the same proof reverts to an ordinary instance of
\cref{prop:task-unavailability}. Sufficiency, throughout this monograph,
is always sufficiency relative to a declared task space, and a proof's
celebrated reliability is precisely a case where that task space has been
closed by convention.

\section{Why Proofs Do Not Exhibit Composition Failure}

Chapter~\ref{ch:budget-relay} onward has treated failure as arising when
a chain's stages pursue \emph{different} local tasks $\tau_0, \ldots,
\tau_{n-1}$, each locally reasonable, that jointly fail to anticipate a
later $\tau^*$. A proof's transmission -- through typesetting, copying,
publication, translation of notation -- is also a chain of renderings in
the sense of \cref{def:substrate-chain}. It is worth stating precisely
why this particular chain typically avoids
\cref{thm:composition-failure}'s failure mode for the verification task,
even though it is exactly the kind of chain the theorem is about in
general.

\begin{proposition}[Homogeneous-Task Chains Preserve What They Share]
\label{prop:homogeneous-chain}
If every stage of a substrate chain shares the same task, $\tau_0 =
\tau_1 = \cdots = \tau_{n-1} = \tau$, and every stage's budget satisfies
$\beta_i \ge R_\tau$ (the total cost of $\tau$'s required distinctions),
then $A_n$ is sufficient for $\tau$.
\end{proposition}

\begin{proofsketch}
By \cref{def:locally-admissible}, each stage ranks $\tau$'s required
distinctions above all others, and by hypothesis $\beta_i \ge R_\tau$ at
every stage, so every required distinction is funded at every stage: no
required distinction is ever dropped, by \cref{def:budget-discard}. Since
this holds for every $i$, $\tau$'s required distinctions survive intact
into $A_n$, and $A_n$ is sufficient for $\tau$ by
\cref{def:sufficiency}.
\end{proofsketch}

\cref{prop:homogeneous-chain}'s hypothesis is exactly what typesetting
and transmission processes for mathematical proofs typically satisfy for
$\tau = \tau_{\text{verify}}$: every party handling the text -- a
typesetter, a journal, a translator of notation -- treats preserving the
logical content as the overriding priority, well-resourced enough that
$\beta_i \ge R_{\tau_{\text{verify}}}$ at every stage. This is why proofs
do not feel like an exception requiring special pleading:
\cref{prop:homogeneous-chain} shows that a chain avoids composition
failure for a given task exactly when every stage treats that task as
priority throughout, which is the normal condition for a proof's
transmission and verification, but conspicuously is \emph{not} the
normal condition for a proof's transmission with respect to explanation.
Nothing about the transmission chain prioritizes preserving the
motivating intuition, the historical context of the discovery, or the
informal picture that made the result plausible in the first place; those
distinctions are exactly the ones a typesetting or publication process
has no particular reason to fund, and \cref{cor:explanation-not-guaranteed}
is the consequence.

\section{Verification as Preservation, Explanation as Reconstitution}

A proof text is high-fidelity (\cref{def:fidelity}) for verification,
since $r_{c,t}$ for a verifier with the text in hand is close to the
proof itself (\cref{def:access-mode}); understanding \emph{why} the
proof works is typically reconstitution-mode, built from a reader's
trace-based reconstruction of the argument's structure, and is therefore
subject to \cref{prop:recursive-insufficiency} in exactly the way memory
is (Chapter~\ref{ch:memory}).
